\documentclass[letterpaper,11pt]{moderncv}
\moderncvstyle{banking}
\moderncvcolor{black}

\usepackage[letterpaper, margin=0.8in]{geometry}
\usepackage{xcolor}
\usepackage{fontspec}
\setmainfont{Times New Roman}
\usepackage[unicode]{hyperref}

\name{Saeyoung}{Kim}
\address{Cambridge, MA, USA}{}{}
\phone[mobile]{+1 (351) 322 5510}
\email{saeyoung\_kim@uml.edu}
\social[linkedin]{saeyoung-macx-kim}

\recipient{Hiring Team}{Boston Dynamics\\Waltham, MA}
\date{\today}
\opening{Dear Hiring Team,}
\closing{Sincerely,}

\begin{document}

\makelettertitle

I'm applying for the Senior Staff Manufacturing Test Engineer role (R2344) on the Atlas program. I've spent the last several years designing and deploying test systems for complex electromechanical hardware, and supporting a research prototype through the transition to reliable, scalable manufacturing is exactly the kind of work I want to be doing.

At Hanwha Systems, I owned the test infrastructure side of military satellite integration across EM, QM, and FM builds. I designed the EGSE architecture, delivered turnkey test setups with data acquisition hardware and sensor interfaces, and wrote the controlled procedures that production teams ran on the factory floor. When electromechanical issues came up during integration, I performed the root cause analysis, isolated whether it was a board-level problem, a wiring fault, or a design gap, and coordinated fixes with the hardware and electrical teams. I pushed design-for-test feedback back into the next build based on failure trends and fault isolation data from each campaign. I also managed purchasing and commissioning of new test stations, trained both engineers and technicians on equipment operation and debug workflows, and worked across electrical, mechanical, software, and program teams to keep test coverage aligned with production schedules.

My PhD was seven years of building test infrastructure from scratch. I designed a multi-modal sensor suite for a CubeSat payload, which meant schematics, PCB layout, firmware in C/C++ over UART, SPI, and I2C, and then building custom test fixtures with oscilloscopes, DMMs, and power supplies to characterize and qualify the hardware. I went through multiple prototype-to-qualification cycles, running environmental and reliability tests each round to find failure modes and close them out. That experience of iterating rapidly through design changes while maintaining test rigor is directly relevant to supporting Atlas as it moves from research into production.

At UMass Lowell, I've been developing Python-based diagnostic software for real-time manufacturing process monitoring, replacing commercial platforms with custom tools for automated anomaly detection across 100+ production batches. This work strengthened my ability to build data-driven test workflows and translate cross-functional requirements into test procedures.

I live in Cambridge, so Waltham is local. Regarding work authorization: I am currently authorized to work in the United States and do not require sponsorship. My green card is pending, and I'm happy to discuss details directly.

\makeletterclosing

\end{document}
